\documentclass[leqno]{article}
\usepackage{verbatim}
\usepackage{array}
\usepackage{listings}
\usepackage{fancyvrb}
\usepackage{enumitem}

\usepackage[utf8]{inputenc}
\usepackage[T1]{fontenc}
\usepackage{textcomp}
\usepackage{multicol} \usepackage{mathtools}
\usepackage{amsmath}
\usepackage{wrapfig}
\usepackage{amssymb}
\usepackage{amsmath,amsfonts,amssymb,amsthm,epsfig,epstopdf,titling,url,array}
\usepackage{hyperref}
\usepackage{eso-pic}
\usepackage{pgf}
\usepackage{tikz}
\usepackage{tikz-cd}
\usepackage{graphicx}

% figure support
\usepackage{import}
\usepackage{xifthen}
\pdfminorversion=7
\usepackage{pdfpages}
\usepackage{transparent}
\usepackage{xcolor}

% geometry
\usepackage{geometry}
\geometry{a4paper, margin=1in}

% paragraph length
\setlength{\parindent}{0em}
\setlength{\parskip}{1em}

\newtheorem*{theorem}{Theorem}
\newtheorem*{lemma}{Lemma}
\newtheorem*{proposition}{Proposition}
\newtheorem*{definition}{Definition}
\newtheorem*{observation}{Observation}

\DeclareMathOperator{\Hom}{Hom}
\DeclareMathOperator{\im}{Im}
\DeclareMathOperator{\coker}{Coker}
\DeclareMathOperator{\coim}{Coim}

\newcommand{\com}[1]{\textcolor{red}{#1}}
\newcommand{\incfig}[1]{%
\center
\def\svgwidth{0.9\columnwidth}
\import{./figures/}{#1.pdf_tex}
}
\newcommand{\incimg}[1]{%
\center
\includegraphics[width=0.9\columnwidth]{images/#1}
}
\pdfsuppresswarningpagegroup=1

\title{Definiciones y Proposiciones GEOALG}
\author{Abel Doñate Muñoz}
\date{}

\begin{document}
\maketitle
\tableofcontents
\newpage

\section{Factores}

\begin{definition}[Resultant] We define the resultant of the polynomials $f(x) = a_0x^n + \cdots + a_n$ and $ g(x) = b_0x^m + \cdots + b_m$ as
   \[
	 \mathcal{R}_x(F, G) = \left| \begin{matrix} 
	   a_0 & \cdots & a_n & 0 & \cdots \\
	   0 & \ddots & \ddots & \ddots & 0\\
	   \vdots & 0 & a_0 & \cdots & a_n \\
	   b_0 & \cdots & b_m & 0 & \cdots \\
	   0 & \ddots & \ddots & \ddots & 0\\
	   \vdots & 0 & b_0 & \cdots & b_m \\
	 \end{matrix}  \right|
  \] 
\end{definition}

\begin{proposition} Let $F(x, y, z), G(x, y, z)$ such that $f(1, 0, 0) \neq 0\neq g(1, 0, 0)$ Then they have a nonconstant homogeneous common factor if and only if $\mathbb{R}_x(F, G) (y, z) =0$
\end{proposition}

\begin{proposition}Properties of the resultant:
  \begin{itemize}[topsep=-6pt, itemsep=0pt]
    \item $\mathbb{R}(F, GH) = \mathbb{R}(F, G) \mathbb{R}(F, H)$
  \end{itemize}
\end{proposition}

\section{Curvas planas en $\mathbb{P}_{\mathbb{C}}^2$}

\com{Espacio proyectivo} $\mathbb{P}^2 = \mathbb{C}^2 \sqcup \mathbb{P}^1 $

En toda la sección $f(x, y, z)$ será un polinomio homogéneo. Definimos la curva 
\begin{definition}[Curva plana] $F=\{(x, y, z): f(x, y, z) \mathbb{P}^2 = 0\}$.
\end{definition}

\begin{proposition}
Relación de Euler. Sea $f$ homogénea de grado $d$, tenemos
 \[
df = xf_x + y f_y + zf_z
\] 
\end{proposition}

\begin{definition}[Singular points] We define the set of singular points as $Sing(F)$.
  \[
	p \in Sing(C) \iff f_x(p) = f_y (p)=f_z(p)=0
  \] 
\end{definition}

\begin{observation} Every projective curve is compact and Hausdorff.
\end{observation}

\begin{proposition} Let $f = \sum_{i=0}^{d} a_ix^iy ^{d-i}$ of degree $d$. Then  $f = \prod _{i=1}^d  (\alpha_i x + \beta _i y)$
\end{proposition}

\begin{definition}[Multiplicity] The multiplicity of $F$ at the point  $p\in F$ is defined as
  \[
	m_p(F) = \min \{m = i+j+k : \frac{\partial^m f}{\partial x^i\partial y^j \partial z^k}(p)\neq 0\}
  \] 
\end{definition}

\begin{definition}[Intersection multiplicity] Let $F, G$ two curves. Let $p = (a, b, c)$
  \[
  i_p(F, G)  = \begin{cases}
	\infty & \text{ if } p \in \text{common component of } F, G\\
	0 & \text{ if } p\neq F\cap G\\
	\max  k : (bz-cy)^k | \mathbb{R}_{x}(F, G)(y, z) & \text{ otherwise}
  \end{cases}
  \] 
  \com{terminar}
\end{definition}

\begin{proposition} Properties of $i_p$ 
  \begin{enumerate}[topsep=-6pt, itemsep=0pt]
    \item $i_p(FG, F) = i_p(F, H) + i_p(G, H)$
  \end{enumerate}
\end{proposition}

\begin{theorem}[Bezout] $F, G$ curves of degree  $n, m$ respectively. Then
   \[
  \sum_{p\in F \cap G} i_p(F, G) = nm
  \] 
\end{theorem}

\begin{proposition} $i_p(F, G)\ge m_p(F)m_p(G)$
\end{proposition}






\end{document}
